% Section content template for individual Educative lessons
% This file contains the actual content for a specific section
% Generated from Educative API section content

% Section: Code for ESPNcricinfo
% Section ID: 5124345484804096
% Section Slug: code-for-espncricinfo
% Generated: 2025-12-14T13:33:28.970521

We've reviewed the various aspects of ESPNcricinfo and examined the attributes associated with the problem using different UML diagrams. Let 's now explore the more practical side of things, where we will work on implementing the ESPNcricinfo problem using multiple languages. This is usually the last step in an object-oriented design interview process.
\\
We have chosen the following languages to write the skeleton code of the different classes present in ESPNcricinfo:

\begin{itemize}
\item
\protect\phantomsection\label{FVUQMGyEq1qx1IFHDDoeV}
Java
\item
\protect\phantomsection\label{BQPC-EPpzSIoE8cqHxJ49}
C\#
\item
\protect\phantomsection\label{jbZ_X10wGhRaR6tuUBN9W}
Python
\item
\protect\phantomsection\label{CVRaDvNZi7CoA_A8fRwyY}
C++
\item
\protect\phantomsection\label{t97jgTbx07QKzAJylKKkh}
JavaScript
\end{itemize}
\\
If you want to skip directly to the implementation and see the complete workflow, simply scroll down or click~\hyperref[Executable-code-Designing-ESPNcricinfo]{\textbf{Executable Code: Designing EspnCricinfo}.}

\subsection{ESPNcricinfo classes}\label{w84oQhMZPtIWnZVcnG0HE}
\\
This section will provide the skeleton code of the classes designed in the class diagram lesson.
\begin{quote}
\textbf{Note:} For simplicity, we are not defining getter and setter functions. The reader can assume that all class attributes are private, accessed through their respective public getter methods, and modified only through their public method functions.
\end{quote}
\subsubsection{Constants}\label{qdNtSWp9zA3fYSCu7ITEz}
\\
The following code defines the various enums and custom data types used in the ESPNcricinfo design:
\begin{quote}
\textbf{Note:} JavaScript does not support enumerations. We'll use the \texttt{Object.freeze()} method as an alternative. It freezes an object, preventing further modifications.
\end{quote}

\begin{lstlisting}[language=Java]
public class Address
{
    private int zipCode;
    private string streetAddress;
    private string city;
    private string state;
    private string country;
}

public enum MatchResult
{
    LIVE,
    BAT_FIRST_WIN,
    FIELD_FIRST_WIN,
    DRAW,
    CANCELED
}

public enum UmpireType
{
    FIELD,
    RESERVED,
    THIRD_UMPIRE
}

public enum WicketType
{
    BOLD,
    CAUGHT,
    STUMPED,
    RUN_OUT,
    LBW,
    RETIRED_HURT,
    HIT_WICKET,
    OBSTRUCTING,
    HANDLING
}

public enum BallType
{
    NORMAL,
    WIDE,
    NO_BALL,
    WICKET
}

public enum RunType
{
    NORMAL,
    FOUR,
    SIX,
    LEG_BYE,
    BYE,
    NO_BALL,
    OVERTHROW
}

public enum PlayingPosition
{
    BATTING,
    BOULING,
    ALL_ROUNDER
}
\end{lstlisting}

\subsubsection{Admin, player, coach, and umpire}\label{jcNNUVoA-pKibSiXmIXyQ}
\\
The definitions of the \texttt{Admin}, \texttt{Player}, \texttt{Coach}, and \texttt{Umpire} classes are as follows:

\textbf{The Admin, Player, Coach, and Umpire classes}

\begin{lstlisting}[language=Java]
public class Admin {
  private List<Player> players = new ArrayList<>();
  private List<Team> teams = new ArrayList<>();
  private List<Match> matches = new ArrayList<>();
  private List<Tournament> tournaments = new ArrayList<>();
  private List<Stat> stats = new ArrayList<>();
  private List<News> newsList = new ArrayList<>();
  
  public boolean addPlayer(Player player);
  public boolean addTeam(Team team);
  public boolean addMatch(Match match);
  public boolean addTournament(Tournament tournament);
  public boolean addStats(Stat stats);
  public boolean addNews(News news);
  
  public boolean assignStadiumToMatch(Stadium stadium, Match match);
  public boolean assignUmpireToMatch(Umpire umpire, Match match);
  public boolean assignCommentatorToMatch(Commentator commentator, Match match);
}

public class Player {
  private String name;
  private int age;
  private int country;
}

public class Coach {
  private String name;
  private int age;
  private int country;
  private List<Team> teams;
}

public class Umpire {
  private String name;
  private int age;
  private int country;

  public boolean assignMatch(Match match);
}
\end{lstlisting}

\subsubsection{Run, ball, wicket, over, and innings}\label{A-al5f_Wb76JjlH6mV8P3}
\\
In cricket, the fundamental concepts can be categorized into five types: run, ball, wicket, over, and innings. To store information about these identities, we defined the \texttt{Run}, \texttt{Ball}, \texttt{Wicket}, \texttt{Over}, and \texttt{Innings} classes as shown below:

\textbf{Implementation of the Run, Ball, Wicket, Over, and Innings classes}

\begin{lstlisting}[language=Java]
public class Run {
  private int totalRuns;
  private RunType type;
  private Player scoredBy;
}

public class Ball {
  private Player balledBy;
  private Player playedBy;
  private BallType type;
  private List<Run> runs;
  private Wicket wicket;

  public boolean addCommentary(Commentary commentary);
}

public class Wicket {
  private WicketType type;
  private Player playerOut;
  private Player balledBy;
  private Player caughtBy;
  private Player runoutBy;
  private Player stumpedBy;
}

public class Over {
  private int number;
  private Player bowler;
  private int totalScore;
  private List<Ball> balls;

  public boolean addBall(Ball ball);
}

public class Innings {
  private Playing11 bowling;
  private Playing11 batting;
  private Date startTime;
  private Date endTime;
  private int totalScores;
  private int totalWickets;
  private List<Over> overs;

  public boolean addOver(Over over);
}
\end{lstlisting}

\subsubsection{Match}\label{IQy0CuY6QMdAUeRSuVRZc}
\\
The \texttt{Match} is an abstract class representing matches in ESPNcricinfo. Matches can be of three types: \texttt{T20}, \texttt{Test}, and \texttt{ODI}. The implementation of these classes is given below:

\textbf{Match and its derived classes}

\begin{lstlisting}[language=Java]
public abstract class Match {
  private Date startTime;
  private MatchResult result;
  private int totalOvers;
  private List<Playing11> teams;
  private List<Innings> innings;
  private Playing11 tossWin;
  private Map<Umpire, UmpireType> umpires;
  private Stadium stadium;
  private List<Commentator> commentators;
  private List<MatchStat> stats;

  public abstract boolean assignStadium(Stadium stadium);
  public abstract boolean assignUmpire(Umpire umpire);
  public boolean addTeam(Playing11 team);
  public boolean assignCommentator(Commentator commentator);
}

public class T20 extends Match {
  public boolean assignStadium(Stadium stadium);
  public boolean assignUmpire(Umpire umpire);
}

public class Test extends Match {
  public boolean assignStadium(Stadium stadium);
  public boolean assignUmpire(Umpire umpire);
}

public class ODI extends Match {
  public boolean assignStadium(Stadium stadium);
  public boolean assignUmpire(Umpire umpire);
}
\end{lstlisting}

\subsubsection{Team, tournament squad, and playing11}\label{ptW5d8lw6M7VpKBmFqwpZ}
\\
The definitions of the \texttt{AdminTeam}, \texttt{TournamentSquad}, and \texttt{Playing11} classes are as follows:

\textbf{The Team, TournamentSquad, and Playing11 classes}

\begin{lstlisting}[language=Java]
public class Team {
  private String name;
  private List<Player> players;
  private Coach coach;
  private List<News> news;
  private TeamStat stats;

  public boolean addSquad(TournamentSquad squad);
  public boolean addPlayer(Player player);
  public boolean addNews(News news);
}

public class TournamentSquad {
  private List<Player> players;
  private Tournament tournament;
  private List<TournamentStat> stats;

  public boolean addPlayer(Player player);
}

public class Playing11 {
  private List<Player> players;

  public boolean addPlayer(Player player);
}
\end{lstlisting}

\subsubsection{Tournament, points table, and stadium}\label{JOi82YXqSnT_Th0_exBV2}
\\
The definitions of the \texttt{Tournament}, \texttt{PointsTable}, and \texttt{Stadium} classes are as follows:

\textbf{The Tournament, PointsTable, and Stadium classes}

\begin{lstlisting}[language=Java]
public class Tournament {
  private Date startDate;
  private List<TournamentSquad> teams;
  private List<Match> matches;
  private PointsTable points;

  public boolean addTeam(TournamentSquad team);
  public boolean addMatch(Match match);
}

public class PointsTable {
  private HashMap<String, float> teamPoints;
  private HashMap<Team, MatchResult> matchResults;
  private Tournament tournament;
  private Date lastUpdated;
}

public class Stadium {
  private String name;
  private Address location;
  private int maxCapacity;
  
  public boolean assignMatch(Match match);
}
\end{lstlisting}

\subsubsection{Commentator, commentary, and news}\label{F7WoLt45BiKA-XrjdBsE1}
\\
The definitions of the \texttt{Commentator}, \texttt{Commentary}, and \texttt{News} classes are as follows:

\textbf{Commentator, Commentary, and News classes}

\begin{lstlisting}[language=Java]
public class Commentator {
  private String name;

  public boolean assignMatch(Match match);
}

public class Commentary {
  private String text;
  private Date createdAt;
  private Commentator commentator;
}

public class News {
  private Date date;
  private String text;
  private List<byte> image;
  private Team team;
}
\end{lstlisting}

\subsubsection{Statistics}\label{9MtpKYbuFpRgBXrGlJ1V6}
\\
The \texttt{Stat} is an abstract class that represents the statistics in ESPNcricinfo. Statistics can be of three types: \texttt{PlayerStat}, \texttt{MatchStat}, and \texttt{TeamStat}. These classes are represented below:

\textbf{Stat and its derived classes}

\begin{lstlisting}[language=Java]
public abstract class Stat {
  public abstract boolean updateStats();
}

public class PlayerStat extends Stat {
  private int ranking;
  private int bestScore;
  private int bestWicketCount;
  private int totalMatchesPlayed;
  private int total100s;
  private int totalHattricks;

  public boolean updateStats();
}

public class MatchStat extends Stat {
  private double winPercentage;
  private Player topBatsman;
  private Player topBowler;

  public boolean updateStats();
}

public class TeamStat extends Stat {
  private int totalSixes;
  private int totalFours;
  private int totalReviews; 

  public boolean updateStats();
}
\end{lstlisting}

\subsection{Executable code: Designing ESPNcricinfo}\label{XcqXZGoAsscstZRETmH8h}
\\
Below is a fully self-contained, runnable program in Java, C\#, C++, Python, and JavaScript that demonstrates the core workflows of the ESPNcricinfo system. The main driver code of the system resides in the~\texttt{Driver.java},~\texttt{Driver.cs},~\texttt{Driver.py},~\texttt{Driver.cpp}, and~\texttt{Driver.js}~for all respective languages. You can click the ``Run'' button to execute the codes.

\subsubsection{What does this code show}\label{-0lvMPhQYkMcoeDw4Hs-0}

\begin{itemize}
\item
\protect\phantomsection\label{bta-pWoTq2xNurEjbQwtU}
\textbf{System initialization:} Initializes the system by creating a new tournament with an associated PointsTable and Admin instance. The tournament is configured with a start date, and its internal lists for teams and matches are prepared. This forms the foundational setup for all cricket-related scenarios.
\item
\protect\phantomsection\label{QGX5PX6hRZTJ_iphu0GZx}
\textbf{Scenario 1 (Match creation and assignment):} An admin creates a T20 match and sets the start time and number of overs. A stadium named ``Wankhede Stadium'' is assigned to the match along with key roles: an umpire (``Aleem Dar'') and a commentator (``Harsha Bhogle''). This scenario demonstrates how a match is structured and populated with essential entities.
\item
\protect\phantomsection\label{5K-6GLDqtsP0Gllgz9e8N}
\textbf{Scenario 2 (Team creation and player assignment):} Two players, Virat Kohli and Mitchell Starc, are created and added to two separate teams using Playing11. These teams are then registered to the match and the tournament as TournamentSquads. This shows how teams and players are formed and tied to both the match and the broader tournament.
\item
\protect\phantomsection\label{uhkaGmVRGa0hLU5cUImb8}
\textbf{Scenario 3 (Ball delivery and commentary):} A ball is simulated where Starc bowls to Kohli, who scores a boundary (4 runs). The ball is linked to a Run object, and a Commentary entry is recorded, highlighting Kohli's stroke. This scenario reflects real-time in-game activity and how gameplay details and commentary are captured.
\item
\protect\phantomsection\label{lZSH5-mzHIs8g6dW6Lcxy}
\textbf{Scenario 4 (Points table update):} The points table is updated after the simulated match: Team 1 earns 2 points, while Team 2 earns none. The latest standings are printed, along with the time stamp of the most recent update. This scenario shows how team performance affects tournament rankings.
\end{itemize}
\begin{quote}
\textbf{Hint: Try customizing the code!}\\
This code is designed not just for reading, but for experimentation. You can~\textbf{edit, extend, or modify}~any part of the example.
\end{quote}

\begin{lstlisting}[language=Java]
import java.util.*;

public class Driver {
    public static void main(String[] args) {
        System.out.println("🏏 === Welcome to Cricinfo System Simulation ===
");

        // ==================================================
        // 🛠️ Initialization: Tournament and Admin Setup
        // ==================================================
        System.out.println("📦 === Initialization ===
");

        Tournament tournament = new Tournament();
        tournament.setStartDate(new Date());
        tournament.setTeams(new ArrayList<>());
        tournament.setMatches(new ArrayList<>());

        PointsTable pointsTable = new PointsTable();
        pointsTable.setTeamPoints(new HashMap<>());
        pointsTable.setMatchResults(new HashMap<>());
        pointsTable.setTournament(tournament);
        pointsTable.setLastUpdated(new Date());
        tournament.setPoints(pointsTable);

        Admin admin = new Admin();

        System.out.println("✅ Tournament created and points table initialized.");
        System.out.println("   ➤ Tournament Start Date: " + tournament.getStartDate());
        System.out.println("   ➤ Points table last updated: " + pointsTable.getLastUpdated() + "
");

        // ==================================================
        // 🧩 Scenario 1: Match Creation and Assignment
        // ==================================================
        System.out.println("🎯 === Scenario 1: Create Match and Assign Roles ===
");

        Match match = admin.createMatch(MatchType.T20);
        match.setStartTime(new Date());
        match.setTotalOvers(20);
        System.out.println("✅ Match Created: T20 Format");
        System.out.println("   ➤ Match Start Time: " + match.getStartTime());
        System.out.println("   ➤ Total Overs: " + match.getTotalOvers());

        Stadium stadium = new Stadium();
        stadium.setName("Wankhede Stadium");
        Address address = new Address();
        address.setCity("Mumbai");
        address.setCountry("India");
        stadium.setLocation(address);
        admin.assignStadiumToMatch(stadium, match);
        System.out.println("✅ Stadium Assigned: " + stadium.getName() + ", " + address.getCity() + ", " + address.getCountry());

        Umpire umpire = new Umpire();
        umpire.setName("Aleem Dar");
        umpire.setAge(50);
        umpire.setCountry(92);
        admin.assignUmpireToMatch(umpire, match);
        System.out.println("✅ Umpire Assigned: " + umpire.getName());

        Commentator commentator = new Commentator();
        commentator.setName("Harsha Bhogle");
        admin.assignCommentatorToMatch(commentator, match);
        System.out.println("✅ Commentator Assigned: " + commentator.getName());

        tournament.addMatch(match);
        System.out.println("📥 Match added to tournament.
");

        // ==================================================
        // 👥 Scenario 2: Add Players and Form Teams
        // ==================================================
        System.out.println("🧢 === Scenario 2: Add Players and Teams ===
");

        Player player1 = new Player();
        player1.setName("Virat Kohli");
        player1.setAge(35);
        player1.setCountry(91);

        Player player2 = new Player();
        player2.setName("Mitchell Starc");
        player2.setAge(34);
        player2.setCountry(61);

        Playing11 team1 = new Playing11();
        team1.addPlayer(player1);

        Playing11 team2 = new Playing11();
        team2.addPlayer(player2);

        match.addTeam(team1);
        match.addTeam(team2);
        System.out.println("✅ Players added to match:");
        System.out.println("   ➤ Team 1: " + player1.getName());
        System.out.println("   ➤ Team 2: " + player2.getName());

        TournamentSquad squad1 = new TournamentSquad();
        squad1.setPlayers(List.of(player1));
        tournament.addTeam(squad1);

        TournamentSquad squad2 = new TournamentSquad();
        squad2.setPlayers(List.of(player2));
        tournament.addTeam(squad2);
        System.out.println("📥 Teams registered to the tournament.
");

        // ==================================================
        // 🏏 Scenario 3: Simulate Gameplay and Commentary
        // ==================================================
        System.out.println("🏃 === Scenario 3: Ball Played and Commentary ===
");

        Ball ball = new Ball();
        ball.setBalledBy(player2);
        ball.setPlayedBy(player1);
        ball.setType(BallType.NORMAL);

        Run run = new Run();
        run.setTotalRuns(4);
        run.setType(RunType.FOUR);
        run.setScoredBy(player1);
        ball.setRuns(List.of(run));

        Commentary commentary = new Commentary();
        commentary.setText("Kohli punches it through the covers for a boundary!");
        commentary.setCreatedAt(new Date());
        commentary.setCommentator(commentator);
        ball.addCommentary(commentary);

        System.out.println("🎬 Ball Delivered:");
        System.out.println("   ➤ Bowler: " + player2.getName());
        System.out.println("   ➤ Batsman: " + player1.getName());
        System.out.println("   ➤ Runs Scored: " + run.getTotalRuns() + " (" + run.getType() + ")");
        System.out.println("   ➤ Commentary: "" + commentary.getText() + "" - " + commentator.getName() + "
");

        // ==================================================
        // 📊 Scenario 4: Update and Display Points Table
        // ==================================================
        System.out.println("📊 === Scenario 4: Points Table Update ===
");

        pointsTable.getTeamPoints().put("Team 1", 4.0f);
        pointsTable.getTeamPoints().put("Team 2", 0.0f);
        pointsTable.setLastUpdated(new Date());

        System.out.println("🏆 Updated Points Table:");
        for (Map.Entry<String, Float> entry : pointsTable.getTeamPoints().entrySet()) {
            System.out.println("   ➤ " + entry.getKey() + ": " + entry.getValue() + " points");
        }

        System.out.println("
✅ Points Table Last Updated On: " + pointsTable.getLastUpdated());
        System.out.println("
🏁 === Cricinfo Tournament Simulation Complete ===");
    }
}
\end{lstlisting}

\subsection{Wrapping up}\label{6jFctsKze0OJS6S6Cfrrq}
\\
We've explored the complete design of ESPNcricinfo in this chapter. We 've looked at how ESPNcricinfo can be visualized using various UML diagrams and designed using object-oriented principles and design patterns.

% End of section content